\documentclass{rvcepaper}
\begin{document}

% ================= BLOCK A: 2025-26 papers (inserted after the Apr-Jun 2026 block)
% ---------- CV266TEE Intelligent Transport Systems, CIE - II, 26-05-2026
\begin{xpaper}
\paperinfo{deptline={Department of Civil Engineering},date={26-05-2026},exam={CIE -- II},maxmarks={60},sem={VI},prog={UG},duration={2 Hrs},
 title={Intelligent Transport Systems},code={CV266TEE}}
\xhead
\ptitle{Quiz}
\begin{xtable}{Questions}
\q{1.}{Differentiate between logical architecture and physical architecture in ITS.}{2}{2}{2}
\q{2.}{State the need for ITS architecture in urban areas.}{2}{1}{2}
\q{3.}{Mention any four communication tools used in ITS.}{2}{1}{3}
\q{4.}{Explain the role of traveller information systems in ITS.}{2}{2}{3}
\q{5.}{List the various detection and identification methods used in ITS.}{2}{1}{3}
\end{xtable}
\ptitle{Test}
\begin{xtable}{Questions}
\q{1.}{Describe equipment packages and market packages in ITS architecture with suitable examples.}{10}{3}{2}
\q{2.}{Explain the ITS architecture with neat sketches of logical architecture and physical architecture.}{10}{3}{2}
\q{3.}{Explain various detection, identification, and data collection methods used in Intelligent Transportation Systems.}{10}{3}{3}
\q{4.}{Explain the development of Traffic Management Systems and the functions of a Traffic Management Centre (TMC).}{10}{3}{3}
\q{5.}{Explain the concepts of Advanced Vehicle Control Systems and Advanced Public Transport Systems.}{10}{3}{3}
\end{xtable}
\end{xpaper}

% ---------- HS251TA Principles of Management and Economics, CIE - II, 15th Dec 2025
\begin{iempaper}
\paperinfo{shownote=0,date={15\textsuperscript{th} Dec 2025},exam={CIE -- II},maxmarks={10 + 50},sem={V},prog={UG},duration={30 + 90 Min},
 title={Principles of Management and Economics},code={HS251TA}}
\iemhead
{\small Note:}\par\vspace{-1pt}
\hspace*{0.5cm}{\footnotesize 1. \hspace{0.1cm}Answer all the Questions.}\par\vspace{3pt}
\begin{cietable}{M}{BT}{CO}
\qpart{Part -- A}
\q{1}{Define organization structure?}{1}{L1}{CO2}
\q{2}{Define the term span of control.}{1}{L2}{CO2}
\q{3}{What does ``threat of substitutes'' refer to in Porters Five force Model?}{1}{L2}{CO2}
\q{4}{Who developed the concept of strategy in 1960s and is also called the father of strategic management?}{1}{L1}{CO2}
\q{5}{What are standing plans?}{1}{L1}{CO2}
\q{6}{If product passes through different stages of production, the suitable form of departmentation is \blank[2cm].}{1}{L2}{CO2}
\q{7}{What is elasticity of demand?}{1}{L1}{CO4}
\q{8}{If the percentage increase in price is less than the percentage decrease in quantity demanded, the good is said to be highly \blank[2cm].}{1}{L2}{CO4}
\q{9}{What do you understand by the term monopolistic competition?}{1}{L2}{CO4}
\q{10}{What does the term Equilibrium in Markets for Goods and Services mean?}{1}{L2}{CO4}
\qpart{Part -- B}
\q{1}{Discuss the key differences between strategic and tactical planning.}{10}{L2}{CO2}
\q{2}{Explain the term corporate strategy. Discuss the importance of corporate strategy for a large scale business enterprise.}{10}{L2}{CO2}
\q{3}{Discuss the effectiveness of mechanistic and organic structures in different environmental conditions.}{10}{L3}{CO2}
\q{4}{Illustrate with examples the various forms of departmentalization.}{10}{L2}{CO2}
\q{5}{What are the characteristic features of an oligopolistic industry? How does it differ from monopolistic competition?}{10}{L2}{CO4}
\end{cietable}
\end{iempaper}

% ================= BLOCK B: older papers (appended at the end of the file)
% ---------- HS251TA Principles of Management and Economics, CIE - II, 07 January 2025
\begin{iempaperB}
\paperinfo{ayear={2024-2025 (ODD Sem)},date={07 January 2025},code={HS251TA},sem={V},exam={CIE -- II},maxmarks={10 + 50},duration={120 Min},
 title={PRINCIPLES OF MANAGEMENT AND ECONOMICS}}
\iemheadB
\begin{cietable}{M}{BT}{CO}
\qpart{Part -- A}
\q{1}{\blank[2cm] plans are specifically designed to meet the needs of a unique situation.}{01}{1}{1}
\q{2}{Creating or changing the organization structure is called \blank[2cm].}{01}{1}{1}
\q{3}{Identify the Quadrant of the BCG Matrix which represents high market share and high market growth.}{01}{1}{1}
\q{4}{\blank[2cm] planning bridges the gap between strategic goals and operational tasks.}{01}{1}{2}
\q{5}{\blank[2cm] refers to how standardized an organization's jobs are and the extent to which employee behaviour is guided by rules and procedures.}{01}{1}{1}
\q{6}{A rigid, tightly controlled structure, characterized by high specialization is a characteristic of \blank[2cm] organization.}{01}{1}{2}
\q{7}{Draw the circular flow diagram with reference to microeconomics.}{02}{1}{4}
\q{8}{The rate of increase in prices over a given period of time is defined as \blank[2cm].}{01}{1}{4}
\q{9}{The \blank[2cm] measures the change in prices paid by consumers for goods and services.}{01}{1}{5}
\qpart{Part -- B}
\q{1}{Many companies have a goal of becoming more environmentally sustainable. One of the most important steps they can take is controlling E-waste. You've been put in charge of creating a program to do this for your company to develop plans. Describe your plans in terms of being: (a) Strategic and operational. (b) Short term and long term.}{10}{2}{2}
\q{2}{Bring out the ways in which a firm could be organized for smooth functioning. How the Chain of command and Span of control does affect the efficient functioning of the firm.}{10}{2}{2}
\q{3}{An organization should be viewed as a unified, purposeful system composed of interrelated parts. Illustrate using food industry as an example.}{10}{4}{1}
\q{4}{Will the firms in an oligopoly act more like a monopoly or more like competitors? Briefly explain with suitable example.}{10}{3}{4}
\q{5}{A rapidly developing country has hired you to present a specific strategy for keeping their inflation and unemployment rates low while maintaining a healthy rate of GDP growth. Explain in detail the components of GDP that can be used to accomplish this.}{10}{3}{4}
\end{cietable}
\btfoot{|l|l|l|*{11}{C{0.8cm}|}}{\hline
\multirow{2}{*}{\shortstack{Marks\\Distribution}} & \multicolumn{2}{c|}{Particulars} & CO1 & CO2 & CO3 & CO4 & CO5 & L1 & L2 & L3 & L4 & L5 & L6\\\cline{2-14}
 & Max & Quiz & 04 & 02 & -- & 03 & 01 & 10 & -- & -- & -- & -- & --\\\cline{3-14}
 & Marks & Test & 10 & 20 & -- & 20 & -- & -- & 20 & 20 & 10 & -- & --\\\hline}
\end{iempaperB}

% ---------- 21IE6F8 Elements of Financial Management, CIE - II, 24th July 2024 (Quiz)
\begin{iempaperB}
\paperinfo{ayear={2023-2024 (EVEN Sem)},dept={INDUSTRIAL ENGINEERING \& MANAGEMENT},date={24\textsuperscript{th} July 2024},code={21IE6F8},sem={VI},exam={CIE -- II},maxmarks={10},duration={30 Min},
 title={Elements of Financial Management (Institutional Elective)}}
\iemheadB
\begin{cietable}{M}{BT}{CO}
\qpart{Part -- A}
\q{1}{In the context of bond valuation, what does the term Yield to maturity refer to?}{1}{2}{CO1}
\q{2}{\blank[2cm] type of bond pays no periodic interest but is issued at a discount.}{1}{1}{CO2}
\q{3}{The model used to determine the expected return on a stock based on its risk relative to the market is \blank[2cm].}{1}{1}{CO2}
\q{4}{Name the valuation method that calculates the present value of a project's cash flows minus its initial investment.}{1}{1}{CO1}
\q{5}{What is the term for the risk associated with an entire market or market segment?}{1}{2}{CO1}
\q{6}{\blank[2cm] refers to a strategy that involves spreading of investments to reduce risk.}{1}{1}{CO1}
\q{7}{What is the term that describes the potential for loss in an investment?}{1}{2}{CO1}
\q{8}{Mention the capital budgeting technique that measures the profitability of a project by calculating the return rate that makes the NPV equal to zero.}{1}{1}{CO2}
\q{9}{The period required for the return on an investment to repay the original investment is \blank[2cm].}{1}{1}{CO2}
\q{10}{\blank[2cm] method evaluates the efficiency of an investment by comparing the accounting profit to the initial investment cost.}{1}{1}{CO2}
\end{cietable}
\btfoot{|l|l|l|*{10}{C{0.85cm}|}}{\hline
\multirow{2}{*}{Marks Distribution} & \multicolumn{2}{c|}{Particulars} & CO1 & CO2 & CO3 & CO4 & L1 & L2 & L3 & L4 & L5 & L6\\\cline{2-13}
 & Quiz & Max Marks & 05 & 05 & -- & -- & 07 & 03 & -- & -- & -- & --\\\hline}
\end{iempaperB}

% ---------- 21IE6F8 Elements of Financial Management, CIE - II, 24th July 2024 (Test)
\begin{iempaperB}
\paperinfo{ayear={2023-2024 (EVEN Sem)},dept={INDUSTRIAL ENGINEERING \& MANAGEMENT},date={24\textsuperscript{th} July 2024},code={21IE6F8},sem={VI},exam={CIE -- II},maxmarks={50},duration={90 Min},
 title={Elements of Financial Management (Institutional Elective)}}
\iemheadB
\begin{cietable}{M}{BT}{CO}
\qpart{Part -- B}
\q{1a}{A Rs.200/-par value bond carries a coupon rate of 12\% and maturity period of 8 years. The interest rate is payable semi-annually. Compute the value of the bond if Required Rate of Return is 14\%.}{5}{3}{CO1}
\q{b}{The expected dividend per share is Rs.12/-. The dividend is expected to grow at the rate of 7\% per year. If the price of the share is Rs.145/- now, determine the Expected rate of Return.}{5}{3}{CO1}
\q{2a}{The current dividend for an equity share is Rs.6/-. If the constant compound growth rate is 6\%, determine the dividend 10 years hence.}{5}{3}{CO1}
\q{b}{The price per bond of M/S. Akshara Industries is Rs.180/-. The par value of the share is Rs.200/-. The coupon rate of the bond is 14\%. The maturity period for the bond is 6 years. Determine the Yield to Maturity for the investor.}{5}{3}{CO1}
\q{3}{The stock of Box Limited performs well relative to other stocks during recessionary periods. The stock of Cox Limited, on the other hand, does well during growth periods. Both the stocks are currently selling for \rupee\,100 per share. You assess the rupee return (dividend plus price) of these stocks for the next year as follows:
\qtabc{|l|c|c|c|c|}{\hline & \multicolumn{4}{c|}{\textbf{Economic Condition}}\\\cline{2-5} & \textbf{High Growth} & \textbf{Low Growth} & \textbf{Stagnation} & \textbf{Recession}\\\hline Probability & 0.3 & 0.4 & 0.2 & 0.1\\\hline Return on Box's Stock & 100 & 110 & 120 & 140\\\hline Return on Cox's Stock & 150 & 130 & 90 & 60\\\hline}
Calculate the expected return and standard deviation of investing:
\begin{rom}\item \rupee\,1,000 in the equity stock of Box Limited\item \rupee\,1,000 in the equity stock of Cox Limited\end{rom}}{10}{4}{CO1}
\q{4a}{Enumerate the steps involved the process of Capital Budgeting.}{5}{2}{CO2}
\q{b}{A project which is being evaluated by a firm has a cost of capital of 12\%. The initial outlay for the project is Rs.2 lakhs. The benefits accrued for the next four years are given below.
\qtabc{|c|c|c|c|c|}{\hline Year & 1 & 2 & 3 & 4\\\hline Amount (Rs) & 50,000 & 80,000 & 80,000 & 1,00,000\\\hline}
Calculate the Benefit Cost Ratio and Net Benefit Cost Ratio for the project. Also based on the standard decision criteria, suggest the feasibility of investing in the project.}{5}{4}{CO2}
\q{5}{Long-term finance is critical for a company's growth and stability, providing the necessary capital to fund substantial investments and strategic initiatives. In the light of the above statement, provide a brief review of different types of sources.}{10}{2}{CO3}
\end{cietable}
\btfoot{|l|l|l|*{10}{C{0.85cm}|}}{\hline
\multirow{2}{*}{Marks Distribution} & \multicolumn{2}{c|}{Particulars} & CO1 & CO2 & CO3 & CO4 & L1 & L2 & L3 & L4 & L5 & L6\\\cline{2-13}
 & Test & Max Marks & 30 & 10 & 10 & -- & -- & 15 & 20 & 15 & -- & --\\\hline}
\end{iempaperB}

\end{document}
